% Isolated proposed B169, checkpoint36. All predecessor bytes are retained.
\subsection*{Full topology of the entire one-sided contact image}
\statement{B.169}{Proved}{The actual one-sided five-generator image has full induced profinite topology, and its whole contact kernel is an exact continuous free family receiving the actual source kernel}
\label{stmt:B169}
Use the entire actual one-sided quotient of B.131.5:
\[
 \Gamma_0\longrightarrow
 G=F(a,B,d,E,f,h,y)\rtimes S,\qquad S=F(m,v).
\]
The four corrected source contacts map to \((a,B,d,E)\).
The subgroup
\[
 H_1=\langle a,d,B,E,m\rangle\le G,\qquad
 H_1\cong\bigl(F(a,d)\times\mathbb Z_m\bigr)*F(B,E)
\]
has its full profinite topology induced by \(G\). Consequently its
actual closure is
\begin{equation}\label{eq:B169-image-completion}
 \mathcal H_1=
   \bigl(\widehat F(a,d)\times\widehat{\mathbb Z}_m\bigr)
       \amalg\widehat F(B,E)\ \le\ \widehat G.
\end{equation}
Here and below \(\amalg\) is the full profinite coproduct.
The kernel of \(\mathcal H_1\to
J_0=\widehat{\mathbb Z}_a\times\widehat{\mathbb Z}_m\)
is the free profinite group on the continuous compact family
\begin{equation}\label{eq:B169-compact-labels}
 d_i\quad(i\in\widehat{\mathbb Z}),\qquad
 B_{i,n},E_{i,n}\quad((i,n)\in\widehat{\mathbb Z}^{\,2}).
\end{equation}
Write this kernel as \(\mathcal U_1\).

In the actual filled source, put
\[
 H_5=\langle A_s,B_s,D_s,E_s,m\rangle,\qquad
 \mathcal H_5=\overline{H_5}\le\widehat{\Gamma_0},\qquad
 \mathcal I_5=\mathcal H_5\cap\ker\widehat q_a,
\]
where \(\widehat q_a\) is the \emph{filled} one-family map of
B.139 and B.154. The whole one-sided quotient gives
\begin{equation}\label{eq:B169-source-triangle}
 \pi_+:\mathcal H_5\twoheadrightarrow\mathcal H_1,\qquad
 \pi_+(\mathcal I_5)=\mathcal U_1,\qquad
 \mathcal I_5\le\mathcal I.
\end{equation}
The group \(\mathcal I\) is the actual known contact-source kernel of
B.164 and B.166. No presentation or embedding of the original
\(H_5\) is asserted.

\begin{proof}
\emph{B.169.1. The entire marked quotient and two cyclic edges.}
The quotient of B.131 kills the four barred family generators in the
entire B.123 presentation. Every signed source operator preserves their
normal closure, and every node loses a barred core. Thus all forty
ordinary seeds, both sheets, all signed translates, and all their normal
conjugates are retained in the justification of this actual quotient.
The free-basis changes and their inverses are
\[
 \begin{gathered}
 E=fde d^{-1}f^{-1},\qquad h=a^{-1}ga,\qquad y=xaf,\\
 e=d^{-1}f^{-1}Efd,\qquad g=aha^{-1},\qquad x=yf^{-1}a^{-1}.
 \end{gathered}
\]
With \(\kappa=f^{-1}h\), the actual \(m\)-action is
\[
 \begin{gathered}
 \alpha(a)=a,\quad\alpha(d)=d,\quad\alpha(y)=y,\qquad
 \alpha(B)=\kappa B\kappa^{-1},\quad
 \alpha(E)=\kappa E\kappa^{-1},\\
 \alpha(f)=f^{-1}hf,\qquad\alpha(h)=f.
 \end{gathered}
\]
Its inverse fixes \(a,d,y\), conjugates \(B,E\) by \(hf^{-1}\),
and sends \(f\) to \(h\), \(h\) to \(hfh^{-1}\).
Put
\[
 Q=F(f,h),\quad G_Q=Q\rtimes\langle m\rangle,\quad
 t=\kappa^{-1}m,\quad E_m=F(a,B,d,E,f,h,y)\rtimes\langle m\rangle.
\]
There is the exact ordinary presentation
\begin{equation}\label{eq:B169-two-edges}
 E_m=G_Q*_{\langle m\rangle}
       \bigl(\langle m\rangle\times F(a,d,y)\bigr)
       *_{\langle t\rangle}
       \bigl(\langle t\rangle\times F(B,E)\bigr).
\end{equation}
Indeed the first new relations are precisely that \(m\) commutes with
\(a,d,y\); the second are \([t,B]=[t,E]=1\), equivalent to the
displayed conjugation formulas. The substitutions in both directions
recover the entire mapping-torus presentation. B.131.5 already proves
the ordinary presentation of \(H_1\).

\emph{B.169.2. The full compact cyclic intersection.}
One has
\[
 \widehat{G_Q}=\widehat F(f,h)\rtimes\widehat{\mathbb Z}_m.
 \tag{B.169.1}
\]
To see cofinality explicitly, given a finite \(G_Q\)-target, take a
characteristic finite-index subgroup of \(Q\) inside its fiber kernel.
The action descends to the finite quotient. Also retain the \(m\)-exponent
modulo the order of its prescribed image. The resulting finite
semidirect target refines the entire original kernel. Conversely these
are actual finite \(G_Q\)-targets. This proves the displayed completion
and the full topology of its fiber.

Let \(M=\overline{\langle m\rangle}\) and
\(T=\overline{\langle t\rangle}\) there. Both are copies of
\(\widehat{\mathbb Z}\), because their source projection is the identity.
Write \(q=hf\), \(q'=fh=fqf^{-1}\). The already proved ordered identities
for every integer \(n\) are
\[
 \kappa_{2n}=q^{-n}(q')^n,\qquad
 \kappa_{2n+1}=q^{-n}(f^{-1}h)(q')^n,\qquad
 t^n=\kappa_n^{-1}m^n.
 \tag{B.169.2}
\]
Their order agrees with
\(\kappa_{n+r}=\alpha^n(\kappa_r)\kappa_n\).
For \(z\in\widehat{\mathbb Z}\), define
\(\kappa_z=m^zt^{-z}\), which lies in the compact fiber.
Continuity extends these identities to every profinite \(n\), separately
on the even and odd clopen cosets. If \(m^z=t^w\), the source projection
gives \(z=w\) and \(\kappa_z=1\). Odd \(z\) is excluded by the fiber
abelianization \((-1,1)\). For \(z=2w\), the equality becomes
\[
 q^{-w}fq^wf^{-1}=1.
\]
The primitive-generator centralizer in the full free profinite group
is \(\widehat{\mathbb Z}_f\), as used in B.153. Thus \(q^w\) lies in that
group. Its \(h\)-exponent is \(w\), forcing \(w=0\). Therefore
\begin{equation}\label{eq:B169-cyclic-intersection}
                         M\cap T=1.
\end{equation}

\emph{B.169.3. Refining every finite target of \(H_1\).}
Let \(\psi:H_1\to K\) be any finite target. Its restriction to
\(\langle m\rangle\) extends to \(M\); write its open kernel as \(M_0\).
Choose \(V\triangleleft_{\mathrm{open}}\widehat{G_Q}\) with
\[
 V\cap M\le M_0,\qquad M\cap TV\le M_0.
 \tag{B.169.3}
\]
For the second choice, the compact sets \(M\cap TV\) decrease to
\(M\cap T=1\); otherwise their intersections with the compact complement
\(M\setminus M_0\) would have the finite-intersection property.
Further refinement gives the first condition.

Put \(D=\widehat{G_Q}/V\) and let \(\overline M,\overline T\) be the
two cyclic images. The prescribed map descends to
\(\rho:\overline M\to K\), with
\(\rho(\overline M\cap\overline T)=1\).
Use the finite induced set
\[
 \Omega=D\times_{\overline M}K,\qquad
 [gh,k]=[g,\rho(h)k]\quad(h\in\overline M),
\]
with \(D\) acting by left multiplication. Its identity fiber
\(\Omega_0=\{[1,k]:k\in K\}\) is a copy of \(K\).
Let \(a,d\) act on this fiber by the left regular permutations
\(\psi(a),\psi(d)\), and trivially on its complement. They commute
with \(\overline M\), since their prescribed images commute with
\(\psi(m)\). Set \(y=1\). For each \(p\in\{B,E\}\), set
\[
 p[t_0,k]=[t_0,\psi(p)k]\qquad(t_0\in\overline T)
\]
on \(\overline T\Omega_0\), and set \(p=1\) on its complement.
If two such representatives define the same fiber, their discrepancy
lies in \(\overline T\cap\overline M\), whose \(\rho\)-image is one.
The definition is therefore well-defined, has the inverse obtained
from \(\psi(p)^{-1}\), and commutes with \(\overline T\).
All relations of \eqref{eq:B169-two-edges} hold.
This gives a finite \(E_m\)-target whose restriction to the invariant
\(\Omega_0\) is the regular representation of \(\psi\).
Its restriction kernel on \(H_1\) lies inside \(\ker\psi\).
Thus \(H_1\) has full topology in \(E_m\).

\emph{B.169.4. Retaining all signed source operators.}
For any finite target \(\eta:E_m\to L\), choose a characteristic
finite-index subgroup \(C\) of the free-seven fiber inside its
\(\eta\)-kernel. All actual signed \(S\)-operators descend to \(P=F_7/C\).
If \(\ell=\operatorname{ord}\eta(m)\), let \(A_S\) be the finite image of
\[
 S\longrightarrow\operatorname{Aut}(P)\times\mathbb Z/\ell,\qquad
 m\longmapsto(\overline\alpha_m,1),\quad
 v\longmapsto(\overline\alpha_v,0).
\]
Then \(G\to P\rtimes A_S\) is an actual finite whole-group map.
If \(w m^n\in E_m\) is killed, then \(w\in C\) and \(\ell\mid n\),
so \(\eta(wm^n)=1\). This refines the entire \(E_m\)-kernel.
The explicit exponent factor is needed; a fiber refinement alone would
not prove this assertion. Combined with the preceding construction,
every finite \(H_1\)-target is refined by an actual finite \(G\)-target.
The completion of its ordinary presentation is
\eqref{eq:B169-image-completion}, by the universal finite-target property
of the displayed commuting and free factors.

\emph{B.169.5. The continuous free family and its inverse maps.}
For every positive integer \(n\), form the free profinite group \(U_n\)
on labels \(d_i\), \(i\in\mathbb Z/n\), and
\(B_{i,j},E_{i,j}\), \((i,j)\in(\mathbb Z/n)^2\).
For divisibility, reduce each label to define the surjective bonding map.
Put \(U=\varprojlim U_n\). Its compact label space is the disjoint union
in \eqref{eq:B169-compact-labels}. A continuous map from that space to a
finite group factors through a finite residue grid, so extends to \(U\).
For a profinite target, apply this construction to every finite quotient;
uniqueness gives compatible homomorphisms. The labels topologically
generate \(U\), since they generate every \(U_n\).
This proves the asserted continuous-family universal property directly.

The action of \(J_0\) is jointly continuous and is
\[
 a^s m^t:d_i\mapsto d_{i+s},\qquad
 B_{i,n}\mapsto B_{i+s,n+t},\qquad
 E_{i,n}\mapsto E_{i+s,n+t}.
\]
It is compatible with the finite-grid actions. In \(U\rtimes J_0\),
the five marked elements satisfy \([a,m]=[d_0,m]=1\), so the presentation
gives a map from \(\mathcal H_1\). In the reverse direction send
\[
 d_i\mapsto a^ida^{-i},\quad
 B_{i,n}\mapsto a^im^nBm^{-n}a^{-i},\quad
 E_{i,n}\mapsto a^im^nEm^{-n}a^{-i}.
\]
These continuous maps intertwine the \(J_0\)-action, using exactly
\([a,m]=[d,m]=1\), and hence define the reverse semidirect-product map.
Both composites fix all five marked generators. Thus
\(\mathcal H_1=U\rtimes J_0\) and \(\mathcal U_1=U\).

\emph{B.169.6. The actual source triangle.}
The continuous one-sided map sends \(\mathcal H_5\) onto
\(\mathcal H_1\), since its compact image contains dense \(H_1\).
By B.154, the filled one-family image of \(\mathcal H_5\) is precisely
the actual \(\widehat{\mathbb Z}^{\,2}\) factor generated by \(A_s,m\).
Its map agrees on the five generators with
\(\mathcal H_5\to\mathcal H_1\to J_0\), and hence agrees everywhere.
Every point of \(\mathcal U_1\) therefore has a preimage in
\(\mathcal I_5\), proving \eqref{eq:B169-source-triangle}.
All these source generators are fixed after the actual node filling,
so \(\mathcal I_5\) lies in the known contact kernel.
The map \(\mathcal H_5\to\mathcal H_1\) is not injective:
B.132 detects \([m,D_s]\ne1\), whereas its image is one.
No equality with the full remaining fixed kernel follows.
\end{proof}
